\section{Constraints and shortlist}\label{sec:constraints}

\paragraph{SAR and peak $\Bone$.} \Cref{fig:pareto}a gives the best joint design (no complex alias)
as a function of the largest nominal angle. Both curves are flat above about $280^\circ$ (worst case)
and $310^\circ$ (mean): larger pulses buy nothing. Below, the price rises steadily. A limit of $240^\circ$
costs $6\%$ in the worst case and $14\%$ in the mean; a limit of $180^\circ$ costs $25\%$ and $52\%$; a
limit of $120^\circ$ more than doubles both. Relative to $330^\circ$, a second pulse at $276^\circ$ needs $30\%$ less
energy at equal pulse duration; at equal peak $\Bone$ it is $16\%$ shorter and needs $16\%$ less energy.

\paragraph{Two-stage pipeline.} When scan 2 serves only to map a smooth $\Bone$, $\alpha_1$ should minimise
the scan-1 $\Tone$ bound with $\Bone$ known (\cref{fig:pareto}b). The minimum lies at $17^\circ$ (worst case,
$109.7$) to $19^\circ$ (mean, $64.5$), against $117.8$ and $68.0$ at $15^\circ$: a gain of $5$--$7\%$, consistent with
the information peak just above the Ernst angles (\cref{fig:theory}c). Between $15^\circ$ and $22^\circ$ the worst case
varies by about $7\%$. Because smoothing makes the noise in the $\Bone$ map negligible~\cite{technote}, $\alpha_2$ is
chosen for an alias-free, well-conditioned $\Bone$ measurement and for robustness; restricted to
$\alpha_2\le330^\circ$ (no more SAR than the published design), the smallest worst-case $\Bone$ bound for
$\alpha_1=17^\circ$ is at $\alpha_2=254^\circ$ (without the restriction it lies near $440^\circ$, a larger lever at $78\%$ more
energy, \cref{fig:maps}c). This pulse crosses $180^\circ$ at $\Bone=0.71$, which costs nothing in the $\Tone$ bound but
produces magnitude aliases for $\Bone$ below $0.8$; $\alpha_2=258$--$270^\circ$ avoids the crossing (e.g.\ $17^\circ/270^\circ$:
worst-case $\Bone$ bound $12.3$ against $11.2$).

\begin{figure}[!htbp]
\centering
\includegraphics[width=\linewidth]{pareto.png}
\caption{(a) Best worst-case and mean joint $\Tone$ bound among designs without a complex alias whose
largest angle does not exceed the abscissa. (b) Two-stage pipeline: scan-1 $\Tone$ bound with $\Bone$ known
against $\alpha_1$ (worst case and mean over $\Bone\in[0.7,1.3]$ and the tissues).}
\label{fig:pareto}
\end{figure}

\paragraph{Shortlist.} \Cref{tab:shortlist} collects the designs carried forward. The rules that
select them are: the published design; the minimum worst-case and minimum mean joint $\Tone$
bound; the best design whose pulses cross no multiple of $180^\circ$ for $\Bone\in[0.7,1.3]$ (it coincides
with the worst-case optimum, $\alpha_2=258$--$276^\circ$); the best such design with $2.5\%$ headroom at both ends
of the range (\emph{margin}, $\alpha_2=264$--$270^\circ$); the two-stage design above; and the best design with
both angles at most $180^\circ$. All are free of complex aliases in the fitting range; the last column
gives the true $\Bone$ for which a magnitude alias exists.

\begin{table}[!htbp]
\centering
\caption{Shortlisted designs. Bounds per unit $\sigma/\Mzero$, all parameters free, worst case (w.) or
mean over $\Bone\in[0.7,1.3]$ and the five tissues; ``par.'': parenchyma only (no CSF). Scan-1 $\Tone$:
two-stage criterion ($\Bone$ known). Ratio bias: largest $|\mathrm d\ln\Tone/\mathrm d\varepsilon|$ over the range. Last column:
true $\Bone$ with a magnitude alias in the fitting range $[0.6,1.4]$ (no design has a complex alias).}
\label{tab:shortlist}
\footnotesize\setlength{\tabcolsep}{3pt}
\begin{tabular}{lccccccccc}
\toprule
& & \multicolumn{3}{c}{$\Tone$ bound} & $\Bone$ & $\Ttwo$ & scan-1 $\Tone$ & ratio & \\
\cmidrule(lr){3-5}
Design & $\alpha_1/\alpha_2$ & w. & mean & w.\ par. & w. & w. & w., $\Bone$ known & bias & magn.\ alias\\
\midrule
published & 15/330 & 103.6 & 50.2 & 68.0 & 19.2 & 117.3 & 117.8 & 2.34 & 0.87--1.30\\
worst case = no crossing & 20/276 & 100.3 & 52.6 & 73.9 & 10.2 & 115.9 & 113.4 & 2.06 & 1.22--1.30\\
mean & 14/309 & 103.6 & 49.2 & 61.9 & 20.5 & 130.1 & 124.3 & 2.25 & 0.99--1.30\\
margin & 21/270 & 101.4 & 53.7 & 63.7 & 10.2 & 114.8 & 115.3 & 2.02 & 0.70--0.72, 1.27--1.30\\
two-stage & 17/254 & 111.4 & 53.6 & 78.5 & 11.2 & 124.0 & 109.7 & 2.04 & 0.70--0.79\\
low SAR & 28/180 & 125.1 & 81.9 & 117.0 & 32.0 & 127.6 & 129.7 & 3.26 & 0.70--1.00, 1.00--1.21\\

\bottomrule
\end{tabular}
\end{table}
